% page 591 of the facsimile (image p-590 of the scan)
% block 170.2 x 246.3 mm ; left margin 31.8 mm
\pgc{10.160mm}{0.000mm}{
\l{\cel{54}583}
\l{}
\l{\sou{tirso}. (\sou{mitol.}) Emblemo di Bacchus, javelino an-cirkondata da}
\l{\cel{3}hedero, pampri, e finanta forme di pino-kono, quan portis}
\l{\cel{3}la bacchus-pastorini. -\cel{2}DEFIRS.}
\l{}
\l{\sou{tisuo}. (\sou{anat.}) Materio koheranta ye qua konsistas la organi di}
\l{\cel{3}la animali, di la vejetanti, asemblajo de fibri interplektita}
\l{\cel{3}od inter-apud-pozita. - EFI.}
\l{}
\l{\sou{titano}. I. (\sou{mitol.}) Singla ek la deaji primitiva, qui preiris}
\l{\cel{3}la Olimpani (t. e. Okeanos, Koyos, Krios, Hperion, e c.) -}
\l{\cel{3}II. (\sou{kemio}) Metaloido tetratoma, di qua la karakteri situesas}
\l{\cel{3}inter siliko e stano. -\cel{1}\sou{Simbolo kemiala}\cel{1}: Ti. - DEFIS.}
\l{}
\l{\sou{titilar}. (\sou{trans}.) Submisar a tre lejera tusheti iterata, sur ula}
\l{\cel{3}parto di la korpo. - EFIS.}
\l{}
\l{\sou{titrar}. I..(\sou{teknol.)}\cel{2}Determinar la proporciono di oro, di arjento,}
\l{\cel{3}unionita a la aloyuro che la metali precoza - la proporciono}
\l{\cel{3}di silko, di lano, en texuro mixura. - II. (kemio). Determinar}
\l{\cel{3}la quanto (fixa) di la reaktivo quan kontenas dissolve ula}
\l{\cel{3}volumino di liquoro. - DEF.}
\l{}
\l{\sou{titulo}. I. Nomo qua expresas distingo pri rango sociala, ofico,}
\l{\cel{3}- II. Nomo qua expresas ofico, grado. - III. Indiko di la temo}
\l{\cel{3}traktata en libro, maxim-multa-kaze kun la nomo di la autoro,}
\l{\cel{3}skribita sur la kovrilo di libro broshita, sur la dorso, sur}
\l{\cel{3}un ek la folii komencala. - IV. Indiko, en la komenco-parto di}
\l{\cel{3}chapitro, di dividuro di libro - di la parto di la temo qua}
\l{\cel{3}traktesas en lu. - DEFRS.}
\l{}
\l{\sou{tizano}. Infuzuro o dekokturo de substanci medikamenta. - DEFIS.}
\l{}
\l{to.\cel{2}La kozo qua esas hike, ma kelke plu fore kam \textquotedbl{}co\textquotedbl{}.(abreviuro}
\l{\cel{5}di \textquotedbl{}ito\textquotedbl{}).}
\l{}
\l{\sou{tofo}. Petro tre poroza, sorto di agreguro kalka, quan onu renkon-}
\l{\cel{3}tras strate en ula tereni, sub la tero vejetantala. - DEFIRS.}
\l{}
\l{\sou{togo}. Kapo-vesto (de drapo, de veluro, de silko, e c.), ronda,}
\l{\cel{3}sen rebordo (o kun rebordo nur tre mikra), kun la suprajo}
\l{\cel{3}plata, plisoza periferie (maxim-multa-kaze). - EFIS.}
\l{}
\l{\sou{tolo.}\cel{1}(\sou{teknol.}) Fero quan onu reduktis a lami per batado o per}
\l{\cel{3}lamino. - FIS.}
\l{}
\l{\sou{toldo}. Kovrilo de dika telo gudrizita, ou ek ledro, qua uzesas}
\l{\cel{3}por shirmar la kargajo di la chareti, di la dilijenci, di la}
\l{\cel{3}navi. -}
\l{}
\l{tolerar. (\sou{trans.)}\cel{2}Permisar tacite, ne impedar, ne interdiktar, to}
\l{\cel{3}quon onu desaprobas, o to per quo onu lezesas, detrimentesas, o}
\l{\cel{3}per quo onu sufras. - DEFIS.}
\l{}
\l{tomo. Singla de la libri, broshura o bindura, ye qui konsista}
\l{\cel{3}verko imprimura. - EFIRS.}
}
